% ============================================================
% SECTION 5: THE BOUNDARY. Written clean 2026-07-26 from exp_AA.
% Evidence: PAPER_RESTRUCTURE/exp_AA_satuluri_regime/{SUMMARY.md,
%   results_armA.csv, results_armB.csv, results_armC.csv,
%   metis_noise.csv, lfr_generation.csv}
% Weight: HEAVY, equal to the whole negative-territory section
% (PAPER_STRUCTURE.md). This carries the paper's novel finding.
% STYLE: no em-dashes, no AI-characteristic phrasing (see STYLE.md).
% ============================================================

Everything to this point has been measured on real networks whose average degree
lies between 5.5 and 32.6, and on that territory the answer is uniformly
negative. That range is not a neutral choice. It is the range in which
sparsification for community detection is normally studied, and it lies entirely
below the one quantitative boundary the founding work
names~\cite{satuluri2011local}: a synthetic sweep in which local similarity
sparsification ``actually outperforms the original clustering starting from
degree 50.'' A negative result confined to one side of a boundary answers half
the question. This section supplies the other half.

\subsection{Design}

Three properties are required of the setting. Average degree must vary over an
order of magnitude while everything else is held fixed. Community labels must be
known rather than inferred. And the granularity artifact that accounts for most
of the apparent gains in this literature must be removable by construction rather
than by control.

Synthetic benchmarks with planted communities supply the first two. A fixed-$k$
partitioner supplies the third. If the number of communities is an input to the
algorithm, then the sparsified and unsparsified arms return the same number of
communities by definition, and no difference between them can be attributed to
one partition being finer than the other.

We generate LFR benchmark graphs~\cite{lancichinetti2008benchmark} with
$n = 10^4$ nodes and degree exponent $2.1$, matching the generator settings of
the original sweep, at nominal average degrees of $10$, $25$, $50$, $100$ and
$200$, at two mixing levels, with three graphs per configuration. Realized degree
and mixing are recorded per graph, and all results are reported against realized
values. Three sparsifiers are applied at matched realized retention: local
similarity sparsification, degree-based sparsification, and uniform random edge
deletion. Detection is performed both by a free-granularity modularity optimizer
and, through the \texttt{pymetis} bindings, by Metis, a fixed-$k$
balance-constrained cut partitioner from the family the original accuracy claim
concerns. Quality is scored on the original graph throughout. Recovery is
measured against the planted labels with chance-corrected indices and a
size-matched random-partition floor.

\subsection{The threshold}

Table~\ref{tab:boundary} reports, for each average degree, the number of
configurations in which similarity-based sparsification improves modularity
measured on the original graph, and the number in which it improves
chance-corrected agreement with the planted communities, out of twelve
configurations per degree.

\begin{table}[t]
\centering
\caption{Fixed-$k$ partitioner (Metis), realized mixing $\approx 0.6$. Counts are
configurations out of twelve in which similarity sparsification improves the
quantity named, with the mean change beside them. Modularity is measured on the
\emph{original} graph. Because $k$ is an input to the partitioner, the sparsified
and unsparsified arms return identical community counts, so no granularity effect
is possible.}
\label{tab:boundary}
\begin{tabular}{@{}lcccc@{}}
\toprule
Avg.\ degree & $\Delta Q > 0$ & mean $\Delta Q$ & $\Delta\mathrm{AMI} > 0$ & mean $\Delta\mathrm{AMI}$\\
\midrule
11.2  & 0/12  & $-0.103$ & 0/12  & $-0.216$\\
24.6  & 0/12  & $-0.046$ & 6/12  & $-0.072$\\
49.4  & 8/12  & $-0.016$ & 9/12  & $+0.042$\\
101.9 & 10/12 & $+0.006$ & 12/12 & $+0.116$\\
220.9 & 12/12 & $+0.009$ & 12/12 & $+0.093$\\
\bottomrule
\end{tabular}
\end{table}

Below average degree $25$, the method is harmful on both measures in every
configuration tested. At $49.4$ it becomes beneficial in two thirds of them. By
$101.9$ it is beneficial in all twelve configurations on recovery and ten of
twelve on the objective. The strongest individual cells occur at aggressive
retention: at average degree $49.4$ and $15\%$ retention, the partition obtained
from the sparsified graph scores $+0.005$ modularity and $+0.152$ AMI above the
partition obtained from the original graph, and at $220.9$ the modularity gain
reaches $+0.012$.

The location of this transition was not chosen by us, and it was not known to us
when the experiment was designed. It coincides with the degree at which the
original authors reported their method beginning to overtake the unsparsified
baseline.

\subsection{Three ways the result could be spurious}

\paragraph{Run-to-run variation}
Metis is not deterministic across option seeds. We repeated the decisive cells
with ten partitioner seeds per graph and compared the \emph{worst} sparsified run
against the \emph{best} unsparsified run. The gain survives that comparison at
every degree at or above $49.4$, with worst-case modularity differences of
$+0.001$, $+0.005$ and $+0.008$, and worst-case AMI differences of $+0.134$,
$+0.102$ and $+0.064$. Seed standard deviations are $0.002$ in modularity and
$0.011$ in AMI, so the effect exceeds the noise it must clear by a factor of
between two and sixty.

\paragraph{The edge budget rather than the selection}
Any sparsifier at $15\%$ retention produces a smaller graph that is denser within
communities, and a fixed-$k$ partitioner may simply behave better on such a
graph. Degree-based and uniform random sparsification, applied to the same graphs
at the same realized retention, do not reproduce the effect at any degree. Their
modularity changes run from $-0.145$ to $-0.003$, and their recovery changes from
$-0.480$ to $+0.032$. What produces the gain is the similarity signal used to
choose which edges to keep, not the number of edges kept.

\paragraph{Granularity}
The partitioner takes the community count as an input, and both arms are run with
the same value, so the two partitions being compared have identical community
counts by construction. This is the artifact that no amount of matching fully
excludes on a free-granularity detector, and here it is excluded structurally. We
also record cluster-size balance and find it slightly improved rather than
degraded under sparsification, so the gain is not purchased by producing more
uneven communities.

\subsection{The same graphs, a different algorithm family}

Running a free-granularity modularity optimizer on the identical graphs gives a
different answer. Modularity gains never exceed the spread of a handful of random
restarts. Recovery gains do appear, and they grow with degree, reaching $+0.096$
AMI, but they do not survive the granularity control. Partitioning the original
graph at a resolution tuned to match the sparsified partition's community count
recovers the planted communities as well or better in twenty-one of the
twenty-nine configurations where the comparison is defined. What appeared to be a
benefit of sparsification is the benefit of a finer partition, obtainable without
touching the graph.

The variable that decides the outcome is therefore the interaction of density
with the detection algorithm's degrees of freedom, not the sparsifier and not the
graph alone. A partitioner constrained to return a fixed number of balanced
groups is helped, above a density threshold, by having its misleading edges
removed. An optimizer free to choose its own granularity already has a cheaper
route to the same improvement, and takes it.

\subsection{Why density matters: the denoising account}

The original authors attributed their gains to the removal of spurious edges,
noting that their strongest results came from an assay network known to contain
many false-positive interactions. That explanation makes a testable prediction.
If sparsification helps by deleting edges that do not belong to the generative
structure, then its benefit should grow with the proportion of such edges.

We inject uniformly random edges into benchmark graphs while holding the planted
labels fixed, and measure recovery against those labels. The prediction holds.
Recovery gain rises monotonically with the injected fraction in every tested
configuration for similarity-based sparsification, with rank correlations of
unity between noise level and gain. At average degree $50$ and $20\%$ retention,
the gain grows from $+0.062$ with no injected noise to $+0.184$ when the number
of spurious edges equals the number of real ones. Degree-based sparsification
shows no such response, which is consistent with its selection signal being a
function of degree rather than of local structure.

Two qualifications apply. The gain does not vanish when no noise is injected, so
denoising adds to an effect that is already present rather than being its sole
cause. And a benchmark graph's noise is random by construction, whereas the
spurious edges produced by a real measurement process need not be. The experiment
establishes the direction of the mechanism, not its magnitude in the field.

\subsection{What the positive result does not say}

Two facts constrain how this finding should be read, and both come from the same
experiments.

Sparsification repairs a weaker detector rather than producing the best available
partition. At every degree tested, a resolution-tuned modularity optimizer
running on the untouched graph recovers the planted communities better than any
sparsified pipeline. At average degree $49.4$ it reaches $0.956$ AMI, against
$0.754$ for the sparsified fixed-$k$ pipeline. Sparsification lifts the fixed-$k$
partitioner from $0.596$ to $0.754$, a real and substantial improvement to that
algorithm, and still leaves it behind an algorithm that never needed
sparsification. The defensible statement is that similarity sparsification
repairs a fixed-$k$ cut partitioner on dense graphs, not that it improves
community detection.

In our implementation the quality gain is also not a speed gain. Computing exact
pairwise similarity costs more than it saves. End to end, the pipeline runs at
$1.49\times$ the baseline at average degree $10$, $1.01\times$ at $100$, and
$0.58\times$ at $200$, which is slower than not sparsifying at all, even though
detection alone accelerates by up to $7.2\times$. The original speedups were
obtained with an approximate similarity computation reported as two orders of
magnitude cheaper, measured against clustering runs that took hours. Both cost
bases differ from ours by enough to account for the discrepancy without any
disagreement about what should be counted.

One part of the original report does not reproduce. At high mixing, where the
planted structure is nearly absent and baseline recovery falls to between $0.03$
and $0.19$, we find no gain under either detector at any degree, whereas the
original sweep reports its largest improvement in that regime. Whether this is a
genuine non-reproduction or a floor effect of the metrics we use cannot be
settled at that mixing level, and we report it unresolved.
